\section{Fields, analytic points, and horizontal ramification}
\label{sec:preliminaries}

Fix the data and the prolonged norm of Theorem~\ref{thm:completion}.
All subsequent field inclusions refer to these choices. We have
\[
 K\subset K_0:=\bar k(X)\subset\widetilde K\subset K^{\mathrm{sep}}
 \subset E,
 \qquad H=\widehat{\widetilde K}\subset E.
\]
The inclusion of the completion $H$ into $E$ is isometric: it identifies
$H$ with the closure of $\widetilde K$ in $E$. Since $k$ is complete, the
restriction of the chosen norm to $\bar k$ is its unique prolongation.
Thus $\Omega:=\widehat{\bar k}$ is a specified subfield of $H$.
It is algebraically closed and has additive value group $\mathbb Q$,
with the normalization $v=-\log_p|\cdot|$. The equality of value groups
before and after completion follows because every nonzero limit has
the same norm as sufficiently close approximants.

We use Berkovich analytifications over $\Omega$. A field embedding
$K_0\hookrightarrow H$ determines a point of $X_{\bar k}(H)$ supported
at the generic point of $X_{\bar k}$. Together with the specified
embedding $\Omega\hookrightarrow H$, it determines a point of
$X_\Omega(H)$, whose analytic image we denote by $\xi$. Its projection
to $X_{\bar k}$ is generic; its support on $X_\Omega$ need not be
generic. In particular, $\xi$ can be an $\Omega$-rational point.
We retain the whole $H$-valued point, including the embedding of its
completed residue field into $H$, whenever we use $\xi$.

If $F/K_0$ is finite, its normalization gives a smooth proper connected
curve over $\bar k$ with a finite map to $X_{\bar k}$. We call the
ramification indices of this map at closed points \emph{horizontal
ramification indices}. The corresponding discrete valuations are trivial
on $\bar k$ and have residue characteristic zero. These indices are
unrelated to possible wild ramification on a characteristic-$p$
reduction model.

Fix an odd prime $\ell\ne p$. Inside $K^{\mathrm{sep}}$, let $L_\ell$
be the compositum of all finite Galois extensions $F/K_0$ whose horizontal
ramification indices divide $\ell$. The unramified extensions are
included. In particular $\widetilde K\subset L_\ell$: the universal
\'etale field over $K_0$ is the same field $\widetilde K$ after constants
are included. One may see this by descending any finite \'etale cover of
$X_{\bar k}$ to a finite extension of $k$ and then composing with the
finite \'etale constant-field cover of $X$.

The field-theoretic goal can be displayed as follows. The top row consists
of the defining inclusions; the two right vertical arrows are the
containments proved in the successive stages.
\[
\begin{CD}
 \widetilde K @>>> L_\ell @>>> K^{\mathrm{sep}}\\
 @VVV @VV{\text{first stage}}V @VV{\text{second stage}}V\\
 H @= H @= H .
\end{CD}
\]
Every arrow is inside $E$. In particular a constructed abstract field
embedding is not enough by itself: for a normal extension $F/K_0$ we
will use normality to identify its image with the specified subfield
$F\subset E$.

We shall also use invariance of finite \'etale covers of a proper variety
under extension of algebraically closed fields
\cite[Expos\'e~X, Corollary~1.8]{SGA1}. Thus a finite \'etale cover of
$Y_\Omega$, for a smooth proper curve $Y/\bar k$, descends up to
isomorphism to a cover of $Y$. The analytic maps built on smaller open
regions need not descend. Their evaluation at the fixed $H$-valued point
will be sufficient.
